\NativeTitle{61}{第七单元素养评价（B）}{（小数加减法）}{40}{4}{一 & 二 & 三 & 四}{ &  &  & }
\NativeSlot{280}{38}{%
\ExamGroup{一、填空。{\normalfont\kaishu\fontsize{10.8bp}{14bp}\selectfont （每空1分，共23分）}}

% p061_s01

\ExamQuestion{1.}{5}%
整数部分最小的计数单位与小数部分最大的计数单位的和是\mbox{（\hspace{1.8em}）}，差是\mbox{（\hspace{1.8em}）}。\par
}
\NativeSlot{355}{36}{%
% p061_s02

\ExamQuestion{2.}{5}%
从\(0.45\)里面去掉\mbox{（\hspace{1.8em}）}个\(0.01\)，结果是\(0.05\)。\par
}
\NativeSlot{380}{36}{%
% p061_s03

\ExamQuestion{3.}{5}%
比\(7.8\)少\(2.81\)的数是\mbox{（\hspace{1.8em}）}，比\(2.36\)多\(1.4\)的数是\mbox{（\hspace{1.8em}）}。\par
}
\NativeSlot{405}{36}{%
% p061_s04

\ExamQuestion{4.}{5}%
把\(7.5\)的小数点向右移动一位后，得到的数比原来多\mbox{（\hspace{1.8em}）}。\par
}
\NativeSlot{431.5}{36}{%
% p061_s05

\ExamQuestion{5.}{5}%
在括号里填上合适的数。\((\quad)-2.4=3.58\)；\allowbreak{}\(0.86+(\quad)=1.16\)；\allowbreak{}\(4-(\quad)=0.64\)。\par
}
\NativeSlot{483}{46}{%
\ExamQuestion{6.}{5}%
在圆圈里填上“＞”“＜”或“＝”。\par\smallskip
\noindent%
% p061_s06
\makebox[.999\linewidth][l]{\(4.3-2.5\ \bigcirc\ 4.3-0.25\)}\par
\noindent%
% p061_s07
\makebox[.999\linewidth][l]{\(3.1+9.84\ \bigcirc\ 3.84+9.1\)}\par
\noindent%
% p061_s08
\makebox[.999\linewidth][l]{\(7.8-(3.8+1.6)\ \bigcirc\ 7.8-(3.8-1.6)\)}\par
}
\NativeSlot{569}{36}{%
% p061_s09

\ExamQuestion{7.}{5}%
在一次跳高测试中，小明跳了\(1.25\)米，比小芳多跳了8厘米，小芳跳了\mbox{（\hspace{1.8em}）}米。\par
}
\NativeSlot{620}{36}{%
% p061_s10

\ExamQuestion{8.}{5}%
一个两位小数加上它的计数单位后是\(1.2\)，这个小数是\mbox{（\hspace{1.8em}）}。\par
}
\NativeSlot{645}{36}{%
% p061_s11

\ExamQuestion{9.}{5}%
甲数是\(10.9\)，比乙数多\(0.75\)，甲、乙两数的和是\mbox{（\hspace{1.8em}）}。\par
}
\NativeSlot{669}{40}{%
% p061_s12

\ExamQuestion{10.}{5}%
按规律填数：\par\noindent （1）\(0.5,1.2,(\quad),2.6,(\quad)\)。\par

% p061_s13

\ExamQuestion{（2）}{5}%
\(5.19,5.16,5.13,(\quad),(\quad),5.04\)。\par
}
\NativeSlot{744}{36}{%
% p061_s14

\ExamQuestion{11.}{5}%
天天家、云云家、学校在同一条直线上。天天家到学校的距离是\(1.36\)千米，云云家到学校的距离是\(0.78\)千米。天天家和云云家最远相距\mbox{（\hspace{1.8em}）}千米，\par
}
